\documentclass[11pt]{article}
\usepackage[margin=1in]{geometry}
\usepackage{amsmath,amssymb,amsthm,mathtools,microtype,hyperref}
\hypersetup{hidelinks}
\newtheorem{theorem}{Theorem}
\newtheorem{lemma}[theorem]{Lemma}
\newcommand{\ddc}{dd^c}
\newcommand{\CP}{\mathbb P}
\newcommand{\C}{\mathbb C}
\newcommand{\FS}{\omega_{\rm FS}}
\title{An incomplete divisorial metric with exact geometric-smoothing data}
\author{}
\date{4 October 2026}
\begin{document}
\maketitle
\begin{abstract}
We construct geometric-smoothing data on the blowup of a product of
projective lines for which the auxiliary divisorial exponential metric
is incomplete on every dense Zariski open positive regular locus.
The initial potential has analytic singularities, its resolved residual
is K\"ahler, and the prescribed source volume is exactly the square of
the background K\"ahler form. All class, discrepancy and smooth-volume
identities are explicit. A rotational cusp on one factor solves the
remaining equation. Its existence follows from ordered barriers and a
radial integral equation. The full-mass solution has finite relative
energy and the Lelong numbers of a minimal potential. Its exceptional
pole is canceled by that minimal divisorial singularity, leaving a
finite-distance boundary. The result concerns the auxiliary metric;
it makes no assertion about completeness of the evolving maximal flow.
\end{abstract}

\section{The compatible completeness question}
We use $\ddc=(i/(2\pi))\partial\bar\partial$.
Let $(X,\omega)$ be compact K\"ahler with an $\omega$-psh initial
potential $\phi_0$ having analytic singularities. On a log resolution
$\pi:Y\to X$ write
\begin{gather}\label{eq:resolved}
 \pi^*\phi_0=\sum_jm_j\log|\sigma_j|^2+\ell_0,\qquad
 \pi^*(\omega+\ddc\phi_0)=\sum_jm_j[D_j]+R_0,
\end{gather}
where $R_0$ is smooth closed semipositive and $D=\sum_jD_j$ is SNC.
The section norms come from smooth Hermitian metrics, whose curvatures
are $\Theta_j$. If $K_Y=\pi^*K_X+\sum_jb_jD_j$, set
\begin{gather}\label{eq:compatible}
 \lambda=\max_j\frac{m_j}{1+b_j},\qquad
 \Theta=\sum_j(1+b_j)\Theta_j,\qquad\theta=\lambda\Theta+R_0,\\
 |\sigma|^2=\prod_j|\sigma_j|^2,\qquad
 \pi^*dV_X=\prod_j|\sigma_j|^{2b_j}dV_Y.\notag
\end{gather}
Here $dV_Y$ is smooth and strictly positive. These are the data in
\cite[Sections 1.3 and 3.1]{DGL}. The class of $\theta$ is big and can
be non-nef. Its designated full-mass finite-energy solution satisfies
\begin{equation}\label{eq:target}
 \langle(\theta+\ddc\rho)^2\rangle
                      =e^\rho|\sigma|^{-2}dV_Y,
\end{equation}
with the density extended by zero on $D$.

\begin{theorem}\label{thm:main}
There are data \eqref{eq:resolved}--\eqref{eq:compatible} on compact
K\"ahler surfaces with $dV_X=\omega^2$, $R_0$ K\"ahler and
$D$ a disjoint union of two smooth curves, for which the solution
of \eqref{eq:target} has full mass and belongs to
$\mathcal E^1(Y,\theta)$ relative to a minimal potential.
Its Lelong numbers equal those of that minimal potential.
Every dense Zariski open set on which its current is smooth and
strictly positive has an incomplete intrinsic metric.
The prescribed pole coefficient is globally smooth and strictly
positive across both components of $D$.
\end{theorem}

\section{A curve cusp with fixed smooth volume}
On $\CP^1$ let $z=Z_1/Z_0$, $r=|z|^2$,
$\alpha=\ddc\log(1+r)$ and $h=r/(1+r)$, so $\int\alpha=1$.
\begin{lemma}\label{lem:cusp}
For every $c>2$ there is a rotational function $u$, smooth on
$\CP^1\setminus\{0\}$, such that
\begin{equation}\label{eq:curve}
 g_1=c\alpha+\ddc u=e^u h^{-1}c\alpha.
\end{equation}
It extends as a $c\alpha$-psh function with zero Lelong number at $0$,
$\int g_1=c$, and $\int|u|\alpha+\int|u|g_1<\infty$.
For $s=-\log r\to\infty$ and a real constant $A$,
\begin{gather}\label{eq:cusp}
 u=\log(2/c)-2\log(s+A)+O(e^{-s}),\\
 u_t=2/(s+A)+O(e^{-s}),\qquad
 u_{tt}=2/(s+A)^2+O(e^{-s}),\qquad t=\log r.\notag
\end{gather}
\end{lemma}
\begin{proof}
Put $L=1-\log h=1+\log(1+1/r)$ and $w=-2\log L$.
Direct differentiation gives
\[
 \ddc w=\left(-\frac2L+\frac2{rL^2}\right)\alpha,\qquad
 Q(r):=\frac{c\alpha+\ddc w}{e^wh^{-1}c\alpha}
       =\frac{crL^2-2rL+2}{c(1+r)}.
\]
Since $L\geq1$ and $c>2$, $Q$ is positive, with positive limits
$2/c$ at $0$ and $1-2/c$ at infinity. Thus
$0<q_-\leq Q\leq q_+<\infty$.
Choose $a_0\leq\min(0,\log q_-)$ and
$b_0\geq\max(0,\log q_+)$. Then $l=w+a_0$, $H=w+b_0$ obey
\[
 c\alpha+\ddc l\geq e^lh^{-1}c\alpha,\qquad
 c\alpha+\ddc H\leq e^Hh^{-1}c\alpha,\qquad l\leq w\leq H.
\]

We give the radial Dirichlet construction explicitly. The exhaustion
$\CP^1\setminus\{|z|\leq\varepsilon\}$ is the disk
$|\zeta|<R=1/\varepsilon$ in $\zeta=1/z$.
With $q=|\zeta|^2$, its Euclidean equation is
\[
 \Delta u=f(q)e^u-g(q),\qquad
 f(q)=\frac{4c}{1+q},\quad g(q)=\frac{4c}{(1+q)^2},
\]
with boundary value $w$. All these radial data are smooth at $q=0$.
For $k>0$, the regular radial solution of
$(-\Delta+k)v=F(q)$ with $V(R^2)=b$ is elementary.
Writing $v=V(q)$ gives $4(qV')'=kV-F$ and the integral equation
\begin{equation}\label{eq:volterra}
 V(q)=A_0+\frac14\int_0^q\log(q/t)\bigl(kV(t)-F(t)\bigr)\,dt.
\end{equation}
Successive integration converges on every finite interval: the iterates
of the positive kernel on $1$ give $q^j/(4^j(j!)^2)$.
The coefficient of $A_0$ is
$\sum_{j\geq0}(kq/4)^j/(j!)^2>0$, so a unique $A_0$
enforces the boundary value. Differentiation gives
\[
 V'(q)=\frac1{4q}\int_0^q(kV-F)\,dt,
\]
with its continuous value at zero. Smooth data imply smooth $V$ by
induction. The usual interior-maximum argument for $-\Delta+k$
proves uniqueness and order preservation of this radial inverse.

Choose $k\geq\sup f e^H$, positive, on the fixed disk.
Let $T(v)$ be that inverse with boundary $w$ and right side
$kv-fe^v+g$. This right side is nondecreasing on $[l,H]$.
The barrier inequalities and the maximum argument give
$l\leq T(l)$, $T(H)\leq H$, and order preservation.
Starting from $l$, the iterates increase and remain below $H$.
Their bounded right sides and the integrated derivative formula give
uniform bounds for their derivatives in $q$ on the fixed disk.
Compactness and monotonicity give uniform convergence; passing in
\eqref{eq:volterra} gives a solution, and the derivative formula
bootstraps it to smoothness, including the center.
For exhaustion radii tending to infinity the same barriers bound
solutions on each fixed disk. The nonlinear identity
\[
 U_q(q)=\frac1{4q}\int_0^q(f(t)e^{U(t)}-g(t))\,dt
\]
gives all derivative bounds there by induction.
Diagonal compactness supplies a smooth solution on the reciprocal
plane, including $z=\infty$, with
$u=-2\log(-\log r)+O(1)$ at $z=0$.

For the sharper expansion the radial equation in $t=\log r$ is
\[
 u_{tt}=\frac{ce^u}{1+e^t}-\frac{ce^t}{(1+e^t)^2}.
\]
Set $v(t)=u(t)+c\log(1+e^t)$ and $V(s)=v(-s)$.
Then
\[
 V''=k_0(s)e^V,\qquad k_0(s)=c(1+e^{-s})^{-(c+1)},
 \qquad V=-2\log s+O(1).
\]
Convexity gives $V'<0$ eventually and $V'\to0$.
The quantity $J=(V')^2-2k_0e^V$ satisfies
$J'=-2k_0'e^V$ and tends to zero. Since
$k_0'=O(e^{-s})$, $e^V\asymp s^{-2}$,
\[
 J(s)=2\int_s^\infty k_0'(\tau)e^{V(\tau)}\,d\tau
                         =O(e^{-s}s^{-2}).
\]
Therefore $(e^{-V/2})'=\sqrt{c/2}(1+O(e^{-s}))$,
and integration gives
$e^{-V/2}=\sqrt{c/2}(s+A)+O(e^{-s})$.
This identity and the ODE prove \eqref{eq:cusp}.
The local subharmonic potential $u+c\log(1+r)$ is bounded above
at the puncture and extends there \cite[Proposition 3.3.25]{JP}.
Its log-log growth gives zero Lelong number. The radial slope tends to
zero, so integration by parts on punctured disks adds no point mass.
Its positive density is asymptotic to
$(i/(2\pi))\,2(r s^2)^{-1}\,dz\wedge d\bar z$.
Distributional Stokes gives mass $c$. The only singular energy integral
is bounded by $C\int^\infty\log s/s^2\,ds$; both asserted integrals
are finite. This proves the lemma.
\end{proof}

\section{Exact compact geometry and initial data}
Take $X=\CP^1_Z\times\CP^1_W$, with normalized FS forms
$\alpha,\beta$, and put
\[
 \omega=6(\alpha+\beta),\qquad dV_X=\omega^2=72\alpha\wedge\beta.
\]
Let $H=\{Z_1=0\}$ and $p=([0:1],[1:0])=(\infty,0)$.
The smooth functions
\[
 h_H=\frac{|Z_1|^2}{|Z_0|^2+|Z_1|^2},\qquad
 \eta=\frac{|Z_0|^2}{|Z_0|^2+|Z_1|^2}
                 +\frac{|W_1|^2}{|W_0|^2+|W_1|^2}
\]
have zero sets $H$ and $\{p\}$ respectively. Define
$\phi_0=2\log h_H+2\log\eta$.
Blow up $p$: $\pi:Y\to X$. Its exceptional curve $E$ is disjoint
from $H_Y=\pi^{-1}H$. This is a smooth compact connected K\"ahler
surface; for instance the projective blowup embeds in a product of
projective spaces. Put $D=H_Y+E$.

Set $|\sigma_H|^2=\pi^*h_H$ and $|\sigma_E|^2=\pi^*\eta$.
The latter really is a smooth Hermitian section norm. Near $p$ write
$x=Z_0/Z_1$, $y=W_1/W_0$. In the blowup chart $x=t$, $y=tv$,
\[
 \pi^*\eta=|t|^2A(t,v),\qquad
 A(t,v)=\frac1{1+|t|^2}+\frac{|v|^2}{1+|tv|^2}>0.
\]
Thus the metric factor extends smoothly and positively across $E$;
the chart with $x,y$ exchanged covers the remaining points.
The factors agree by the Hermitian transition rule.

The sections $Z_0W_0,\sqrt2Z_0W_1,Z_1W_1$ of $\mathcal O(1,1)$
have summed squared norm $\eta$ and base ideal $(x,\sqrt2xy,y)=(x,y)$.
After division by $\sigma_E$ they have no common zero and define
$F:Y\to\CP^2$. If $\gamma$ is its normalized FS form, the induced
metric gives the exact identity
\begin{equation}\label{eq:curvatures}
 F^*\gamma=\pi^*(\alpha+\beta)-\Theta_E,\qquad
 \Theta_H=\pi^*\alpha.
\end{equation}
In the first blowup chart $F=[1:\sqrt2tv:v]$, so it is an
isomorphism from $E$ onto a line.
The initial logarithms have analytic singularities: near $p$,
\[
 \eta=\frac{|x|^2+|y|^2+2|xy|^2}{(1+|x|^2)(1+|y|^2)},
\]
and near $H$ the other term is smooth. The integer coefficient $2$
can also be expressed with the usual reciprocal-integer coefficient
by taking holomorphic powers in the sum of squares.
Globally $\pi^*\phi_0=2\log|\sigma_H|^2+2\log|\sigma_E|^2$,
so $\ell_0=0$, $m_H=m_E=2$, and
\begin{align}\label{eq:residual}
 \pi^*(\omega+\ddc\phi_0)&=2[H_Y]+2[E]+R_0,\\
 R_0&=\pi^*\omega-2\Theta_H-2\Theta_E
      =2\pi^*\alpha+4\pi^*\beta+2F^*\gamma.\notag
\end{align}
The residual is K\"ahler: away from $E$ its first two terms are
positive; on $E$ their kernel is $TE$ and $F^*\gamma$ is positive
there. This also proves that $\phi_0$ is $\omega$-psh.

The Jacobian of $(t,v)\mapsto(t,tv)$ is $t$, giving
$b_E=1$, $b_H=0$. Integrability at $H$ requires an exponential
coefficient $q<1/2$, whereas at the isolated point it requires $q<1$.
Thus $\lambda=2$ directly and from \eqref{eq:compatible}.
Consequently
\begin{equation}\label{eq:theta}
 \theta=2\Theta_H+4\Theta_E+R_0=\pi^*\omega+2\Theta_E,
 \qquad\{\theta\}=\pi^*\{\omega\}+2\{E\}.
\end{equation}
Its intersection with $E$ is $-2$, so the class is non-nef.
It is big by the explicit K\"ahler current
\[
 \theta+\ddc(3\log|\sigma_E|^2)
           =(\pi^*\omega-\Theta_E)+3[E],\qquad
 \pi^*\omega-\Theta_E=5\pi^*(\alpha+\beta)+F^*\gamma>0.
\]
Finally define $dV_Y=\pi^*\omega^2/|\sigma_E|^2$ off $E$.
The blowup Jacobian contributes $|t|^2$ in the numerator, canceled
exactly by $|t|^2A(t,v)$ in the denominator. It therefore extends
smoothly and strictly positively everywhere and satisfies
\begin{equation}\label{eq:volume}
 \pi^*dV_X=|\sigma_E|^2dV_Y.
\end{equation}
The unweighted pole $dV_Y/(|\sigma_H|^2|\sigma_E|^2)$ is
nonintegrable along each component, with smooth strictly positive
coefficient. All identities in the flow setup have now been checked.

\section{The exact solution and its minimal reference}
Take $u$ from Lemma~\ref{lem:cusp} with $c=6$ on the first factor,
and write $g_X=(6\alpha+\ddc u)+6\beta$. Then
\[
 g_X^2=e^uh_H^{-1}\omega^2.
\]
It is smooth and strictly positive near $p$. Put
\begin{equation}\label{eq:rho}
 P=2\log|\sigma_E|^2,\qquad\rho=P+\pi^*u.
\end{equation}
Poincar\'e--Lelong and \eqref{eq:theta} give
\begin{equation}\label{eq:current}
 \theta+\ddc\rho=2[E]+\pi^*g_X.
\end{equation}
Off $D$, the required density is identically
\[
 e^\rho(|\sigma_H|^2|\sigma_E|^2)^{-1}dV_Y
       =e^{\pi^*u}|\sigma_H|^{-2}|\sigma_E|^2dV_Y
       =\pi^*g_X^2.
\]
Both measures extend by zero on $D$, so this is the global
non-pluripolar equation \eqref{eq:target}.
The prescribed unweighted pole is unchanged. Its exceptional factor
is suppressed by the solution's minimal divisorial singularity.

We verify minimality without assuming it from the class description.
For any positive closed current $T$ in $\pi^*\{\omega\}+2\{E\}$,
let $S=\pi_*T$. It is positive closed in $\{\omega\}$, and its
pullback by local psh potentials is well defined and positive.
Off $E$, $T=\pi^*S$. Thus $T-\pi^*S$ is closed, has order zero and
is supported on $E$. The support theorem
\cite[Chapter III, Corollary 2.14]{Dem} makes it a constant multiple
of $[E]$; cohomology fixes that constant to $2$, since $E^2=-1$.
By the K\"ahler $\partial\bar\partial$ lemma write
$S=\omega+\ddc v$ with $v$ $\omega$-psh. Every $\theta$-psh
potential is consequently $P+\pi^*v$ plus a constant.
Every such $v$ is bounded above on compact $X$. Hence $P$ has
minimal singularities.

Its minimal current is $2[E]+\pi^*\omega$, whose non-pluripolar
square has mass $\int_X\omega^2=72$. This is the volume of the
big class \cite[Definition 1.17]{BEGZ}; the ordinary cup-product square
is $68$ and is not the relevant mass. Lemma~\ref{lem:cusp} gives
\[
 \int_Y\langle(\theta+\ddc\rho)^2\rangle
                         =\int_Xg_X^2=72.
\]
Thus $\rho$ has full mass. The difference $\rho-P=\pi^*u$
is integrable for this measure, since $g_X^2=12g_1\wedge\beta$ and
$\int|u|g_1<\infty$. Shift the minimal reference by a constant so that $\rho$ lies below it.
The full-mass energy criterion
\cite[Proposition 2.11(ii)]{BEGZ} proves relative $\mathcal E^1$
membership. In fact all mixed energy integrals are finite:
\[
 \omega^2=72\alpha\wedge\beta,\qquad
 g_X\wedge\omega=6g_1\wedge\beta+36\alpha\wedge\beta,
 \qquad g_X^2=12g_1\wedge\beta,
\]
and $\int|u|\alpha$ is also finite. These mixed products put no
mass on the two divisors. Near $E$, $\pi^*u$ is smooth, giving
exactly the Lelong coefficient $2$ of $P$. The cusp on $H_Y$ has
zero Lelong number; elsewhere everything is smooth. This verifies
minimal Lelong numbers everywhere. For completeness, if $v$ is another
full-mass solution, the comparison principle \cite[Corollary 2.3]{BEGZ}
on $\{v<\rho\}$ gives
$\int e^\rho|\sigma|^{-2}dV_Y\leq\int e^v|\sigma|^{-2}dV_Y$
over that set. Both integrals are finite. Strict inequality of the
integrands forces that set to have zero smooth volume; interchange
$v$ and $\rho$ to obtain equality almost everywhere. Quasi-psh
representatives then agree everywhere. Thus \eqref{eq:rho} is the
equation's designated solution.

\section{An unavoidable finite-distance boundary}
The positive smooth locus of \eqref{eq:current} is
$U_0=Y\setminus(H_Y\cup E)$. The cusp density diverges on $H_Y$;
the nonzero integration current $2[E]$ excludes all points of $E$.
Even the smooth metric part $\pi^*g_X$ is degenerate on $TE$.
Let $U$ be any dense Zariski open positive regular locus.
Its proper analytic complement has finitely many curve and point
components \cite[Theorem 7.8 and Corollary 7.10]{Adamus} and contains
$E$. Choose an exceptional point avoiding its intersections with all
other components. A small coordinate neighborhood there has complement
exactly $E$. The transverse path $(t,v_0)$, $0<t\leq\varepsilon$,
lies in $U$ and tends to that removed point.
Since $g_X$ is smooth near $p$ and $\pi$ is smooth, $\pi^*g_X$ has
bounded coefficients in the whole neighborhood. This path has bounded
speed and finite length. Its tail lengths give an intrinsic Cauchy
sequence with no limit in $U$, since a positive smooth Riemannian
metric induces the manifold topology. Every such $U$ is incomplete,
proving Theorem~\ref{thm:main}.

The example verifies the complete compatible data for the auxiliary
class $\lambda\Theta+R_0$, including $dV_X=\omega^2$.
It does not identify this auxiliary metric with the evolving maximal
flow at other times. The distinct nefness branch has a separate
counterexample with zero Lelong numbers; here non-nefness allows the
minimal exceptional component that cancels its prescribed pole. Completeness
under the simultaneous imposition of nefness and these compatible-data
identities is a stronger question left open by the two examples.

\begin{thebibliography}{9}
\bibitem{Adamus} J.~Adamus, \emph{Complex Analytic Geometry},
lecture notes, revised 8 December 2015,
\href{https://www.math.uwo.ca/faculty/adamus/teaching/9607b2019/CAGI.pdf}{author's notes}.
\bibitem{BEGZ} S.~Boucksom, P.~Eyssidieux, V.~Guedj and A.~Zeriahi,
Monge--Amp\`ere equations in big cohomology classes,
\emph{Acta Mathematica} 205 (2010), 199--262,
\href{https://arxiv.org/abs/0812.3674}{arXiv:0812.3674}.
\bibitem{Dem} J.-P.~Demailly, \emph{Complex Analytic and Differential Geometry},
\href{https://www-fourier.univ-grenoble-alpes.fr/~demailly/manuscripts/agbook.pdf}{author's manuscript}.
\bibitem{DGL} E.~Di Nezza, V.~Guedj and C.~H.~Lu,
Geometric smoothing by the K\"ahler--Ricci flow,
\href{https://arxiv.org/abs/2607.08325}{arXiv:2607.08325v1},
Sections 1.3 and 3.1, Theorem 3.4 and Question 3.5.
\bibitem{JP} M.~Jarnicki and P.~Pflug, \emph{Separately analytic functions},
manuscript, 11 November 2008,
\href{https://www.math.stonybrook.edu/~ebedford/PapersForM537/JarnickiPflugSepHolo081111.pdf}{source copy}.
\end{thebibliography}
\end{document}
